\section*{Abstract}
The growth of digital financial services has made cross-institutional fraud faster to
spread than any single bank can detect, yet privacy regulation and competitive distrust
prevent institutions from pooling raw transaction data. This thesis presents FL-ADTCN, a
blockchain-integrated federated learning framework for collaborative financial fraud
detection that resolves the privacy, trust, and incentive barriers simultaneously. Each
institution trains a local Adaptive Deep Temporal Context Network (ADTCN) -- a temporal
one-dimensional convolutional detector -- on its own transactions, and only model
parameters, never raw data, are shared. The federated aggregation layer combines
\emph{differentially private} weight sharing (Gaussian mechanism), \emph{Krum} robust
selection to reject outlying updates, and \emph{Shapley-value} contribution attribution that
quantifies each institution's marginal value to the global model; the Shapley weights are
written to a Hyperledger Fabric ledger and drive an on-chain, chaincode-enforced token
incentive. The central contribution is a \textbf{characterisation of the
privacy--incentive coupling} of this mechanism: we show that differential privacy and
contribution-based incentives are in direct tension, and that applying the privacy noise on
the released contribution channel rather than on the shared weights restores rank-faithful,
DP-protected on-chain rewards at roughly two orders of magnitude lower privacy budget
($\epsilon^\star$: $3000\!\to\!50$). This is supported by a characterisation suite that bounds
the mechanism: the privacy--utility decoupling of accuracy versus reward fidelity; the economic
isolation of a colluding majority; statistical Byzantine fault tolerance; and the cost and
fidelity limits of Shapley attribution as the federation grows. The supporting substrate is a
Dynamic Butterfly--Billiards Optimisation Algorithm (DB-BOA) that automates ADTCN
hyperparameter tuning (without manual search, though it does not exceed a hand-tuned default) and
consensus leader selection, together with a deliberately minimal reinforcement-learning (linear
Q-learning) policy for \emph{sequential} leader rotation that, in simulation, matches the
optimiser on reward while adapting to nodes whose reliability degrades at runtime. The framework
is implemented end-to-end on a real Hyperledger Fabric consortium with the \texttt{DBBOAContract}
smart contract and evaluated on the public ULB Credit Card Fraud dataset (284{,}807 transactions,
0.17\% fraud) partitioned across three banks. Detection metrics are reported at the deployment
operating point, in which weight-channel differential privacy is disabled and privacy is instead
enforced on the incentive channel, the configuration the contribution above shows to be both
private and useful. \emph{Scalability} is established for contribution \emph{attribution} -- a
Monte-Carlo Shapley estimator that remains tractable to $n=20$ organisations, faithful to
$n\approx 5$ -- rather than for blockchain throughput or distributed training, which remain in the
limitations. \emph{Security} is evaluated on two complementary fronts: Krum is run in the regime
where its $n\ge 2f+3$ guarantee holds ($n=5,f=1$ and $n=7,f=2$), rejecting four weight-level
poisoning attacks, while the economic incentive layer isolates a behavioural colluding majority
that out-numbers any statistical defence; the default three-organisation pipeline runs at $f=0$
(consensus alignment, no adversary assumed).

\vspace{1cm}
\textbf{Keywords: }Federated Learning; FL-ADTCN; Adaptive Deep Temporal Context Network;
Hyperledger Fabric; Permissioned Consortium Blockchain; Smart Contracts; Financial Fraud
Detection; Differential Privacy; Krum; Shapley Value; Contribution Attribution; Token
Incentive Mechanism; Dynamic Butterfly--Billiards Optimisation (DB-BOA); Leader Selection;
Privacy Preservation; Byzantine Resilience
\pagebreak
